\Needspace{32\baselineskip}
\section{Focus insertion, erasure, and recovery}
\label{sec:erasure}

We now erase the focus of an arbitrary focussed graph with no outputs and
prove that reinsertion recovers the original graph and focus up to isomorphism.
Unlike the previous section's recovery result, the input need not have been
constructed by our insertion function. The converse round trip, insertion
followed by erasure, recovers the underlying graph. These constructions
support the state-search semantics developed in Section~\ref{sec:state-search}.

\Needspace{16\baselineskip}
\subsection{Contraction maps and the retained carriers}
Fix a focussed graph $G:I\to\epsilon$, focus edge $e_f$, source $s$, and an
internal vertex $t$ equal to the focus target. We first work with this explicit
witness, then prove that every such graph supplies it.

\begin{definition}[General focus erasure maps]\label{term:focus-erasure}
Remove $t$ from the internal carrier and $e_f$ from the edge carrier, using
subtypes of names unequal to the deleted names. Let $j$ include the remaining
ports into the old ports. If $w\ne t$, let $k(w)$ keep $w$ with its non-deletion
certificate. Self-acyclicity ensures $s\ne t$, so $p=k(s)$ is defined. Set
\[
 c(w)=\begin{cases}p&w=t,\\k(w)&w\ne t,\end{cases}
 \qquad
 a(x)=\begin{cases}t&j(x)=s,\\j(x)&j(x)\ne s.\end{cases}
\]
The map $c$ contracts the focus. The map $a$ recovers the appropriate old
\emph{source} representative, choosing $t$ at the joined port; $j$ recovers
the old \emph{target} representative. Boundary ports and their order are
unchanged.
\end{definition}
\lean{Graph/FocusErasure}{11}{11}{Linear.FocusErasure.Vertex}
\lean{Graph/FocusErasure}{13}{13}{Linear.FocusErasure.lift}
\lean{Graph/FocusErasure}{15}{19}{Linear.FocusErasure.keep}
\lean{Graph/FocusErasure}{38}{38}{Linear.FocusErasure.position}
\lean{Graph/FocusErasure}{40}{42}{Linear.FocusErasure.contract}
\lean{Graph/FocusErasure}{62}{64}{Linear.FocusErasure.expand}

\begin{lemma}[Port-map identities]
The map $j$ is injective, avoids $t$, and satisfies $j(k(w))=w$.
Moreover $j(p)=s$, $c(t)=c(s)=p$, $j(c(w))=w$ for $w\ne t$, and $c(j(x))=x$.
The expansion $a(x)$ never equals $s$; it satisfies $a(c(w))=w$ for $w\ne s$
and $c(a(x))=x$.
\end{lemma}
\begin{proof}
The inclusion and keeping claims follow by port-constructor cases; subtype
extensionality gives injectivity on internal ports, and the subtype certificate
excludes $t$. The inequality $s\ne t$ is exactly focus self-acyclicity rewritten
by the chosen target witness. The identities for $p,c$ follow by unfolding and
splitting on whether the argument is $t$. Apply inclusion injectivity to derive
$c(j(x))=x$ from $j(c(j(x)))=j(x)$.
For expansion, when $j(x)=s$ its result is $t\ne s$; otherwise it is already
different from $s$. To prove $a(c(w))=w$, split on $w=t$. At $t$, contraction
lands at $p$ and expansion returns $t$. Otherwise inclusion after contraction
returns $w\ne s$, so expansion leaves it unchanged. For $c(a(x))=x$, split
on $j(x)=s$. In that case $x=p$ by inclusion injectivity; otherwise use
$c(j(x))=x$.
\end{proof}
\lean{Graph/FocusErasure}{21}{22}{Linear.FocusErasure.lift_keep}
\lean{Graph/FocusErasure}{24}{25}{Linear.FocusErasure.lift_injective}
\lean{Graph/FocusErasure}{27}{31}{Linear.FocusErasure.lift_ne_target}
\lean{Graph/FocusErasure}{33}{35}{Linear.FocusErasure.source_ne_target}
\lean{Graph/FocusErasure}{44}{45}{Linear.FocusErasure.lift_position}
\lean{Graph/FocusErasure}{47}{48}{Linear.FocusErasure.contract_target}
\lean{Graph/FocusErasure}{50}{52}{Linear.FocusErasure.lift_contract}
\lean{Graph/FocusErasure}{54}{56}{Linear.FocusErasure.contract_lift}
\lean{Graph/FocusErasure}{58}{59}{Linear.FocusErasure.contract_source}
\lean{Graph/FocusErasure}{66}{71}{Linear.FocusErasure.expand_ne_source}
\lean{Graph/FocusErasure}{73}{78}{Linear.FocusErasure.expand_contract}
\lean{Graph/FocusErasure}{80}{86}{Linear.FocusErasure.contract_expand}

\begin{lemma}[Injectivity on retained incidence lists]
A source of a retained edge cannot equal $s$, and a target of a retained edge
cannot equal $t$. Contraction is injective on vertices other than $s$, and also
on vertices other than $t$.
\end{lemma}
\begin{proof}
If a retained edge sourced $s$, source uniqueness with the focus edge would
identify the two edges, contrary to retention. Target uniqueness proves the
other exclusion. For injectivity off $s$, apply $a$ to equality of contracted
ports and use $a(c(w))=w$. For injectivity off $t$, apply $j$ and use
$j(c(w))=w$. These are restricted injectivity assertions: $c$ itself identifies
the two distinct focus endpoints.
\end{proof}
\lean{Graph/FocusErasure}{88}{91}{Linear.FocusErasure.source_avoids}
\lean{Graph/FocusErasure}{93}{97}{Linear.FocusErasure.target_avoids}
\lean{Graph/FocusErasure}{99}{102}{Linear.FocusErasure.source_injective}
\lean{Graph/FocusErasure}{104}{107}{Linear.FocusErasure.target_injective}

\begin{lemma}[Label preservation]
Both focus endpoints have term label. Keeping an undeleted port preserves its
label, and contraction preserves every port's label.
\end{lemma}
\begin{proof}
Rewrite the focus source and target typing certificates using their singleton
incidences and focus label; its signature is term-to-term. Keeping preserves
labels by cases on internal and boundary ports. For contraction, all ports
other than $t$ use keeping. At $t$, its image is the kept source and both
endpoints have the same term label.
\end{proof}
\lean{Graph/FocusErasure}{109}{112}{Linear.FocusErasure.source_term}
\lean{Graph/FocusErasure}{114}{118}{Linear.FocusErasure.target_term}
\lean{Graph/FocusErasure}{120}{122}{Linear.FocusErasure.keep_label}
\lean{Graph/FocusErasure}{124}{132}{Linear.FocusErasure.contract_label}

\Needspace{16\baselineskip}
\subsection{A valid focus-free underlying graph}
\begin{theorem}[General erasure construction]
Delete $e_f$ and $t$, retain every other edge label, and map all retained source
and target lists through $c$. The result is a finite typed hypernet of the same
interface $I\to\epsilon$.
\end{theorem}
\begin{proof}
Finite subtype enumerations supply the carriers. Restricted contraction
injectivity on retained sources and targets preserves no-repetition. Given a
shared contracted source of two retained edges, restricted injectivity recovers
a shared old source, so old source uniqueness identifies the edge names.
The target proof is the same with the target restriction.

For source coverage, there are no output ports. Given an erased port $x$, use
the old source coverage at $a(x)$. Since $a(x)\ne s$, its source edge cannot be
the focus; contraction of that source gives $c(a(x))=x$.
For target coverage, a retained target $w$ differs from $t$, so
$j(c(w))=w$ and its input status is unchanged. Conversely, if $x$ is not an
input, old target coverage at $j(x)$ gives an edge. It is not the focus because
$j(x)\ne t$. Its contracted target is $c(j(x))=x$.
Finally, mapping contracted lists through labels equals reading the old labels,
by label preservation, so both signature-typing certificates follow from $G$.
\end{proof}
\lean{Graph/FocusErasure}{135}{210}{Linear.FocusErasure.erase}

\begin{lemma}[Erased invariants]
The erased graph is globally focus-free, its distinguished port $p$ has term
label, and erasure preserves global hole-freedom.
\end{lemma}
\begin{proof}
Every retained edge was different from the unique focus, so its entire label,
including any box body, is focus-free. The port $p$ keeps the term-labelled
source, hence has term label. A hypothetical hole in the erased graph lies at
a retained edge with unchanged label and label-local position; the same
occurrence would be a hole in the original graph. Thus a hole-free original
has a hole-free erasure.
\end{proof}
\lean{Graph/FocusErasure}{212}{213}{Linear.FocusErasure.erased_focusFree}
\lean{Graph/FocusErasure}{215}{217}{Linear.FocusErasure.position_term}
\lean{Graph/FocusErasure}{219}{222}{Linear.FocusErasure.erased_holeFree}

\Needspace{16\baselineskip}
\subsection{Reinsertion recovers an arbitrary focussed graph}
\begin{definition}[Restoring a deleted name]\label{term:focus-round-trip}
For any type $A$ and distinguished $b\in A$, the restoration bijection
$(A\setminus\{b\})+1\to A$ sends retained names to themselves and the fresh
singleton name to $b$. Its inverse tests whether a name equals $b$.
Applying this to $t$ and $e_f$ gives the vertex and edge bijections needed for
reinsertion. Reinsert the original focus mode at $p$ in the erased graph.
The port $p$ is a non-output because the output interface is empty, and its
term-typing certificate was proved above.
\end{definition}
\lean{Graph/FocusRoundTrip}{12}{23}{Linear.FocusErasure.restore}
\lean{Graph/FocusRoundTrip}{25}{29}{Linear.FocusErasure.position_not_output}
\lean{Graph/FocusRoundTrip}{32}{34}{Linear.FocusErasure.reinsert}
The bijection's inverse laws follow by retained/fresh cases; the fresh carrier
has exactly one member. At ports, restoring host inclusion agrees with $j$,
and restoring the fresh insertion vertex gives $t$. Restoring the insertion's
source-redirection map agrees with $a$: the test $x=p$ is equivalent to
$j(x)=s$, by inclusion injectivity and $j(p)=s$.
\lean{Graph/FocusRoundTrip}{36}{38}{Linear.FocusErasure.restore_host}
\lean{Graph/FocusRoundTrip}{40}{41}{Linear.FocusErasure.restore_fresh}
\lean{Graph/FocusRoundTrip}{43}{55}{Linear.FocusErasure.restore_split}

\begin{theorem}[Focus-preserving round trip]
Reinserting after erasure is isomorphic to the original graph. The isomorphism
restores the focus edge, source, and target, preserves boundary order, and
preserves literal edge labels, including complete boxed bodies.
\end{theorem}
\begin{proof}
Use the restoration bijections. Retained internal labels are unchanged; the
fresh insertion vertex has term label, equal to the original label of $t$.
Retained edge labels are unchanged, and the new focus label equals the original
focus label by its certificate.
For a retained source list, restoration after insertion-redirection after
contraction acts as $a\circ c$, which is the identity on these sources because
they avoid $s$. For a retained target list, restoration acts as $j\circ c$,
which is the identity because these targets avoid $t$. Map these pointwise
identities over the lists. The new focus source restores to $j(p)=s$, and its
target restores to $t$, giving the original singleton incidence lists.
The fresh edge restores to $e_f$ by definition, and the preceding two port
identities explicitly identify its endpoints.
\end{proof}
\lean{Graph/FocusRoundTrip}{59}{104}{Linear.FocusErasure.reinsert_iso}
\lean{Graph/FocusRoundTrip}{107}{107}{Linear.FocusErasure.reinsert_focus_edge}
\lean{Graph/FocusRoundTrip}{110}{116}{Linear.FocusErasure.reinsert_focus_endpoints}

\begin{lemma}[Every closed focussed graph admits erasure]
The focus target of any graph $G:I\to\epsilon$ is internal. Therefore every
such focussed graph has a globally focus-free underlying graph, preserving
hole-freedom, and reinserting its mode recovers it up to the proved isomorphism.
\end{lemma}
\begin{proof}
The focus target occurs in a target list, so cannot be an input. There are no
outputs. The remaining port constructor is internal. Choose its internal name
with its equality certificate, and instantiate the preceding erasure,
invariant, and round-trip constructions. This use of choice supplies a witness
already proved to exist; it does not assume an erasure law.
\end{proof}
\lean{Graph/FocusRoundTrip}{127}{133}{Linear.Focussed.target_internal}
\lean{Graph/FocusRoundTrip}{135}{135}{Linear.Focussed.targetVertex}
\lean{Graph/FocusRoundTrip}{137}{137}{Linear.Focussed.targetVertex_spec}
\lean{Graph/FocusRoundTrip}{140}{141}{Linear.Focussed.underlying}
\lean{Graph/FocusRoundTrip}{143}{144}{Linear.Focussed.underlying_focusFree}
\lean{Graph/FocusRoundTrip}{146}{147}{Linear.Focussed.underlying_holeFree}
\lean{Graph/FocusRoundTrip}{151}{153}{Linear.Focussed.underlying_reinsert_iso}
The theorem proves the erasure-then-insertion direction for arbitrary closed
focussed graphs. The opposite insertion-then-erasure isomorphism is proved in
Section~\ref{sec:insertion-cancellation}, and naturality under source
isomorphisms in Section~\ref{sec:erasure-renaming}.

\Needspace{16\baselineskip}
\subsection{Focus-preserving isomorphisms and state recovery}
\begin{definition}[Focussed isomorphism]\label{term:focussed-isomorphism}
A focussed isomorphism is a literal-label graph isomorphism together with
equality of the two focus modes and an equation mapping the distinguished
focus edge to the distinguished focus edge. Identity, inverse, and composition
are obtained from the corresponding graph isomorphisms, using the edge
bijection's inverse laws and composition for the focus equations.
\end{definition}
\lean{Graph/StateSearch}{12}{15}{Linear.Focussed.Iso}
\lean{Graph/StateSearch}{20}{20}{Linear.Focussed.Iso.refl}
\lean{Graph/StateSearch}{22}{27}{Linear.Focussed.Iso.symm}
\lean{Graph/StateSearch}{29}{34}{Linear.Focussed.Iso.trans}

\begin{lemma}[Endpoints and observations are invariant]
A focussed isomorphism maps both focus endpoints correctly and preserves
initiality and finality.
\end{lemma}
\begin{proof}
Specialize the isomorphism's source-list equation to the focus edge and rewrite
both lists as singletons. Singleton injectivity identifies the endpoints.
Repeat for target lists. Port maps preserve input status; together with the
mode equality, this yields both boundary-predicate equivalences.
\end{proof}
\lean{Graph/StateSearch}{36}{39}{Linear.Focussed.Iso.source}
\lean{Graph/StateSearch}{41}{44}{Linear.Focussed.Iso.target}
\lean{Graph/StateSearch}{46}{49}{Linear.Focussed.Iso.initial_iff}
\lean{Graph/StateSearch}{51}{54}{Linear.Focussed.Iso.final_iff}

\begin{definition}[Extracted state coordinate]\label{term:state-coordinate}
For an arbitrary state $s$, erase its focus and take the typed position whose
mode is the old mode and whose vertex is the joined port $p$. Its typing is
the erased-position typing lemma. Re-realizing this coordinate in the globally
focus-free, hole-free underlying graph yields a state focus-preservingly
isomorphic to $s$, using the general round trip. Thus the input need not be a
canonical insertion or carry an assumed representation certificate.
\end{definition}
\lean{Graph/StateSearch}{63}{65}{Linear.State.coordinate}
\lean{Graph/StateSearch}{69}{74}{Linear.State.recover}
The extracted coordinate is initial or final exactly when $s$ is. Compose the
previous realization boundary equivalence with the focussed isomorphism's
boundary equivalence.
\lean{Graph/StateSearch}{76}{81}{Linear.State.coordinate_boundary}


\Needspace{20\baselineskip}
\subsection{Insertion followed by erasure}
\label{sec:insertion-cancellation}
Fix a focus-free graph $G:I\to[]$, a term vertex $v$, and a mode $m$.
Write $S=\operatorname{insert}(G,v,m)$ for the focussed graph constructed
in Section~\ref{sec:realization}. The preceding round trip started with an
arbitrary focussed graph. We now establish the opposite direction, including
the position at which the focus was inserted.

\begin{definition}[Deleting a fresh singleton]
For any carrier $A$, deleting the fresh member from $A+\mathbf1$ leaves a
carrier canonically bijective with $A$. The forward map sends a retained
$\mathsf{inl}(a)$ to $a$; the inverse sends $a$ to that retained element.
An element of the right summand cannot survive, since its singleton index
must be zero. Apply general erasure to the fresh focus and fresh target of
$S$ to define $E(S)$.
\end{definition}
The inverse laws hold by cases on the summand: the retained case is reflexive,
and the fresh case contradicts its non-deletion certificate.
\lean{Graph/InsertionRoundTrip}{10}{21}{Linear.FocusInsertion.deleteFresh}
\lean{Graph/InsertionRoundTrip}{27}{28}{Linear.FocusInsertion.erasedInsertion}

\begin{lemma}[Cancellation on ports]
Let $D$ be the port map induced by deletion of the fresh singleton, $j$ the
inclusion of retained ports, and $\kappa$ the collapse map of insertion.
Then $D(x)=\kappa(j(x))$. The joined vertex maps to $v$, and for every port
$w$ of $S$, $D(c(w))=\kappa(w)$, where $c$ is general focus contraction.
\end{lemma}
\begin{proof}
The first identity follows by cases on the retained port. An internal
right-summand case is impossible; input ports are fixed. For the joined
vertex, inclusion gives the focus source, namely the hosted copy of $v$;
collapse sends it to $v$. Finally, if $w$ is the fresh target, contraction
gives the joined vertex and both sides equal $v$. Otherwise inclusion after
contraction is $w$, and the first identity applies.
\end{proof}
\lean{Graph/InsertionRoundTrip}{31}{40}{Linear.FocusInsertion.deleteFresh_port}
\lean{Graph/InsertionRoundTrip}{43}{47}{Linear.FocusInsertion.erase_position}
\lean{Graph/InsertionRoundTrip}{50}{59}{Linear.FocusInsertion.erase_contract}

\begin{theorem}[Insertion--erasure cancellation]
\termdef{insertion-cancellation}
There is an isomorphism $E(S)\cong G$ preserving full labels and ordered
incidences. The choice-based operation $S.\mathsf{underlying}$ has such an
isomorphism too, and one can choose it to send the joined vertex to $v$.
\end{theorem}
\begin{proof}
Use deletion of the singleton on both internal vertices and edges. All
surviving labels are inherited. Sources were first split by insertion and
then contracted by erasure; the port cancellation identity reduces their
composite to collapse after split, which is the identity. Targets were
hosted by insertion, so the analogous composite is collapse after hosting,
also the identity. These pointwise identities give equality of ordered lists.

The target selected by general erasure is necessarily the fresh internal
vertex: both it and the fresh vertex equal the focus target, and the internal
port constructor is injective. Substituting this equality identifies the
general underlying graph with $E(S)$. For the position-sensitive assertion,
carry an arbitrary target together with its equality certificate, identify
it with the fresh target, and use the explicit isomorphism and joined-port
identity together. This keeps graph and coordinate recovery synchronized.
\end{proof}
\lean{Graph/InsertionRoundTrip}{62}{100}{Linear.FocusInsertion.erase_iso}
\lean{Graph/InsertionRoundTrip}{103}{104}{Linear.FocusInsertion.targetVertex_eq}
\lean{Graph/InsertionRoundTrip}{107}{111}{Linear.FocusInsertion.underlying_iso}
\lean{Graph/InsertionRoundTrip}{115}{129}{Linear.FocusInsertion.underlying_position_exists}
